\documentclass[12pt,a4paper]{article}
\usepackage{../../../myconfig}

\begin{document}

% --- KHUNG TIÊU ĐỀ ĐẦU TRANG ---
\titlebox{Chủ đề: Oxyz}{Phương trình mặt cầu}{Oxyz: Phương trình mặt cầu}

% =========================================================================
% TRANG 1: TỪ CÂU 1 ĐẾN CÂU 7 (NỀN TẢNG, DẠNG 1, DẠNG 2, DẠNG 3)
% =========================================================================

% Câu 1: Nền tảng - Xác định tâm và bán kính từ dạng chính tắc
\questionLinesManual{0}{1}{%
Trong không gian $Oxyz$, cho mặt cầu $(S)\colon (x-1)^2+(y+2)^2+(z-3)^2=16$. Tọa độ tâm $I$ và bán kính $R$ của mặt cầu $(S)$ là:
}{%
\begin{tasks}(2)
    \task $I(-1;2;-3), R=4$
    \task $I(1;-2;3), R=4$
    \task $I(1;-2;3), R=16$
    \task $I(-1;-2;-3), R=4$
\end{tasks}
}

% Câu 2: Nền tảng - Xác định tâm và bán kính từ dạng khai triển tổng quát
\questionLinesManual{0}{2}{%
Trong không gian $Oxyz$, cho mặt cầu $(S)\colon x^2+y^2+z^2-4x+2y-6z-2=0$. Bán kính của mặt cầu $(S)$ bằng:
}{%
\begin{tasks}(4)
    \task $R=4$
    \task $R=\sqrt{14}$
    \task $R=16$
    \task $R=2\sqrt{3}$
\end{tasks}
}

% Câu 3: Nền tảng - Điều kiện tồn tại phương trình mặt cầu
\questionLinesManual{0}{3}{%
Trong không gian $Oxyz$, tìm tất cả các giá trị của tham số $m$ để phương trình $x^2+y^2+z^2-2x-2y-4z+m=0$ là phương trình của một mặt cầu.
}{%
\begin{tasks}(4)
    \task $m < 6$
    \task $m \ge 6$
    \task $m \le 6$
    \task $m > 6$
\end{tasks}
}

% Câu 4: Dạng 1 - Viết phương trình mặt cầu biết tâm và bán kính
\questionLinesManual{0}{4}{%
Trong không gian $Oxyz$, phương trình của mặt cầu $(S)$ có tâm $I(0;-2;1)$ và bán kính $R=5$ là:
}{%
\begin{tasks}(2)
    \task $x^2+(y+2)^2+(z-1)^2=25$
    \task $x^2+(y-2)^2+(z+1)^2=25$
    \task $x^2+(y+2)^2+(z-1)^2=5$
    \task $x^2+(y-2)^2+(z+1)^2=5$
\end{tasks}
}

% Câu 5: Dạng 2 - Mặt cầu có tâm I và đi qua điểm A
\questionLinesManual{0}{5}{%
Trong không gian $Oxyz$, cho hai điểm $I(1;1;1)$ và $A(1;2;3)$. Phương trình mặt cầu có tâm $I$ và đi qua điểm $A$ là:
}{%
\begin{tasks}(2)
    \task $(x-1)^2+(y-1)^2+(z-1)^2=5$
    \task $(x+1)^2+(y+1)^2+(z+1)^2=5$
    \task $(x-1)^2+(y-1)^2+(z-1)^2=25$
    \task $(x+1)^2+(y+1)^2+(z+1)^2=29$
\end{tasks}
}

% Câu 6: Dạng 3 - Mặt cầu nhận đoạn thẳng AB làm đường kính
\questionLinesManual{0}{6}{%
Trong không gian $Oxyz$, cho hai điểm $A(2;4;1)$ và $B(-2;2;-3)$. Phương trình của mặt cầu có đường kính $AB$ là:
}{%
\begin{tasks}(2)
    \task $x^2+(y-3)^2+(z-1)^2=36$
    \task $x^2+(y-3)^2+(z+1)^2=9$
    \task $x^2+(y+3)^2+(z-1)^2=9$
    \task $x^2+(y-3)^2+(z+1)^2=36$
\end{tasks}
}

% Câu 7: Dạng 3 - Biến thể đường kính MN
\questionLinesManual{0}{7}{%
Trong không gian $Oxyz$, cho hai điểm $M(3;-2;5)$ và $N(-1;6;-3)$. Mặt cầu đường kính $MN$ có phương trình là:
}{%
\begin{tasks}(2)
    \task $(x+1)^2+(y+2)^2+(z+1)^2=36$
    \task $(x-1)^2+(y-2)^2+(z-1)^2=6$
    \task $(x-1)^2+(y-2)^2+(z-1)^2=36$
    \task $(x+1)^2+(y+2)^2+(z+1)^2=6$
\end{tasks}
}

\newpage

% =========================================================================
% TRANG 2: TỪ CÂU 8 ĐẾN CÂU 14 (DẠNG 4, DẠNG 5, DẠNG 6, DẠNG 7)
% =========================================================================

% Câu 8: Dạng 4 - Tâm I tiếp xúc mặt phẳng tọa độ (Oxz)
\questionLinesManual{0}{8}{%
Trong không gian $Oxyz$, mặt cầu có tâm $I(-1;-2;3)$ và tiếp xúc với mặt phẳng tọa độ $(Oxz)$ có phương trình là:
}{%
\begin{tasks}(2)
    \task $(x+1)^2+(y+2)^2+(z-3)^2=4$
    \task $(x+1)^2+(y+2)^2+(z-3)^2=9$
    \task $(x+1)^2+(y+2)^2+(z-3)^2=1$
    \task $x^2+y^2+z^2+2x+4y-6z+12=0$
\end{tasks}
}

% Câu 9: Dạng 4 - Tâm I tiếp xúc trục tọa độ Oy
\questionLinesManual{0}{9}{%
Trong không gian $Oxyz$, cho điểm $I(2;1;-3)$. Mặt cầu tâm $I$ và tiếp xúc với trục $Oy$ có phương trình là:
}{%
\begin{tasks}(2)
    \task $(x-2)^2+(y-1)^2+(z+3)^2=4$
    \task $(x-2)^2+(y-1)^2+(z+3)^2=13$
    \task $(x-2)^2+(y-1)^2+(z+3)^2=9$
    \task $(x-2)^2+(y-1)^2+(z+3)^2=10$
\end{tasks}
}

% Câu 10: Dạng 5 - Tâm I tiếp xúc với mặt phẳng (P) (ĐÃ SỬA CHUẨN ĐÁP ÁN A)
\questionLinesManual{0}{10}{%
Trong không gian $Oxyz$, cho điểm $I(2;1;-2)$ và mặt phẳng $(P)\colon 2x-2y+z+5=0$. Mặt cầu $(S)$ có tâm $I$ và tiếp xúc với mặt phẳng $(P)$ có phương trình là:
}{%
\begin{tasks}(2)
    \task $(x-2)^2+(y-1)^2+(z+2)^2=\dfrac{25}{9}$
    \task $(x-2)^2+(y-1)^2+(z+2)^2=9$
    \task $(x+2)^2+(y+1)^2+(z-2)^2=9$
    \task $(x-2)^2+(y-1)^2+(z+2)^2=3$
\end{tasks}
}

% Câu 11: Dạng 5 - Tâm I tiếp xúc với mặt phẳng tổng quát
\questionLinesManual{0}{11}{%
Trong không gian $Oxyz$, cho điểm $M(1;2;-3)$ và mặt phẳng $(P)\colon x+2y-2z-2=0$. Viết phương trình mặt cầu tâm $M$ và tiếp xúc với mặt phẳng $(P)$.
}{%
\begin{tasks}(2)
    \task $(x-1)^2+(y-2)^2+(z+3)^2=9$
    \task $(x+1)^2+(y+2)^2+(z-3)^2=9$
    \task $(x-1)^2+(y-2)^2+(z+3)^2=81$
    \task $(x-1)^2+(y-2)^2+(z+3)^2=3$
\end{tasks}
}

% Câu 12: Dạng 6 - Mặt cầu đi qua bốn điểm O, A, B, C
\questionLinesManual{0}{12}{%
Trong không gian $Oxyz$, mặt cầu ngoại tiếp tứ diện $OABC$ với $O(0;0;0)$, $A(1;0;0)$, $B(0;-2;0)$ và $C(0;0;4)$ có phương trình là:
}{%
\begin{tasks}(2)
    \task $x^2+y^2+z^2+x-2y+4z=0$
    \task $x^2+y^2+z^2-x+2y-4z=0$
    \task $x^2+y^2+z^2-2x+4y-8z=0$
    \task $x^2+y^2+z^2+2x-4y+8z=0$
\end{tasks}
}

% Câu 13: Dạng 6 - Ngoại tiếp tứ diện tổng quát
\questionLinesManual{0}{13}{%
Trong không gian $Oxyz$, cho bốn điểm $A(1;0;0)$, $B(0;1;0)$, $C(0;0;1)$, $D(1;1;1)$. Bán kính mặt cầu ngoại tiếp tứ diện $ABCD$ bằng:
}{%
\begin{tasks}(4)
    \task $\sqrt{2}$
    \task $\dfrac{\sqrt{3}}{2}$
    \task $\sqrt{3}$
    \task $\dfrac{3}{4}$
\end{tasks}
}

% Câu 14: Dạng 7 - Qua 3 điểm và tâm thuộc mặt phẳng (P)
\questionLinesManual{0}{14}{%
Trong không gian $Oxyz$, cho ba điểm $M(2;3;3)$, $N(2;-1;-1)$, $P(-2;-1;3)$. Mặt cầu đi qua ba điểm $M, N, P$ và có tâm thuộc mặt phẳng $(\alpha)\colon 2x+3y-z+2=0$ có phương trình là:
}{%
\begin{tasks}(2)
    \task $x^2+y^2+z^2+4x-2y+6z+2=0$
    \task $x^2+y^2+z^2-2x+2y-2z-10=0$
    \task $x^2+y^2+z^2-4x+2y-6z-2=0$
    \task $x^2+y^2+z^2-2x+2y-2z-2=0$
\end{tasks}
}
\drawdashlinesfull

% =========================================================================
% TRANG 3: TỪ CÂU 15 ĐẾN CÂU 20 (DẠNG 8, VỊ TRÍ ĐIỂM & VỊ TRÍ 2 MẶT CẦU)
% =========================================================================

% Câu 15: Dạng 8 - Cắt mặt phẳng theo giao tuyến đường tròn
\questionLinesManual{0}{15}{%
Trong không gian $Oxyz$, cho điểm $I(1;2;-1)$ và mặt phẳng $(P)\colon x-2y-2z-8=0$. Phương trình mặt cầu tâm $I$ cắt mặt phẳng $(P)$ theo một đường tròn có bán kính $r=4$ là:
}{%
\begin{tasks}(2)
    \task $(x+1)^2+(y+2)^2+(z-1)^2=5$
    \task $(x-1)^2+(y-2)^2+(z+1)^2=9$
    \task $(x-1)^2+(y-2)^2+(z+1)^2=25$
    \task $(x+1)^2+(y+2)^2+(z-1)^2=25$
\end{tasks}
}

% Câu 16: Dạng 8 - Biến thể thiết diện đường tròn theo diện tích S
\questionLinesManual{0}{16}{%
Trong không gian $Oxyz$, cho mặt phẳng $(\alpha)\colon 2x+y-z-1=0$ và mặt cầu $(S)\colon (x-1)^2+(y-1)^2+(z+2)^2=4$. Mặt phẳng $(\alpha)$ cắt $(S)$ theo đường tròn có diện tích bằng:
}{%
\begin{tasks}(4)
    \task $\dfrac{16\pi}{3}$
    \task $\dfrac{4\pi}{3}$
    \task $\dfrac{2\pi\sqrt{3}}{3}$
    \task $4\pi$
\end{tasks}
}

% Câu 17: Chuyên đề vị trí tương đối giữa điểm và mặt cầu (ĐÃ ĐỔI ĐIỂM M ĐỂ ĐÁP ÁN B ĐÚNG)
\questionLinesManual{0}{17}{%
Trong không gian $Oxyz$, cho mặt cầu $(S)\colon (x-1)^2+(y+2)^2+(z-3)^2=25$ và điểm $M(m;1;3)$. Tìm tất cả các giá trị thực của tham số $m$ để điểm $M$ nằm ngoài mặt cầu $(S)$.
}{%
\begin{tasks}(2)
    \task $-3 < m < 5$
    \task $m < -3$ hoặc $m > 5$
    \task $m \le -3$ hoặc $m \ge 5$
    \task $-5 < m < 3$
\end{tasks}
}

% Câu 18: Chuyên đề vị trí tương đối giữa 2 mặt cầu - Ở ngoài nhau
\questionLinesManual{0}{18}{%
Trong không gian $Oxyz$, cho hai mặt cầu $(S_1)\colon (x-1)^2+(y-2)^2+(z-3)^2=4$ và $(S_2)\colon (x-4)^2+(y-6)^2+(z-3)^2=4$. Vị trí tương đối giữa $(S_1)$ và $(S_2)$ là:
}{%
\begin{tasks}(2)
    \task Tiếp xúc ngoài nhau
    \task Ở ngoài nhau
    \task Cắt nhau theo đường tròn
    \task Tiếp xúc trong nhau
\end{tasks}
}

% Câu 19: Chuyên đề vị trí tương đối giữa 2 mặt cầu - Tiếp xúc trong
\questionLinesManual{0}{19}{%
Trong không gian $Oxyz$, cho hai mặt cầu $(S_1)\colon x^2+y^2+z^2=25$ và $(S_2)\colon (x-1)^2+(y-2)^2+(z-2)^2=4$. Khẳng định nào sau đây đúng?
}{%
\begin{tasks}(2)
    \task $(S_1)$ và $(S_2)$ tiếp xúc ngoài
    \task $(S_1)$ và $(S_2)$ cắt nhau
    \task $(S_1)$ và $(S_2)$ tiếp xúc trong
    \task $(S_2)$ đựng trong $(S_1)$
\end{tasks}
}

% Câu 20: Chuyên đề vị trí tương đối giữa 2 mặt cầu - Đựng nhau (trong nhau)
\questionLinesManual{0}{20}{%
Trong không gian $Oxyz$, xét vị trí tương đối giữa hai mặt cầu $(S_1)\colon (x-1)^2+(y+2)^2+(z-1)^2=36$ và $(S_2)\colon (x-2)^2+(y+1)^2+z^2=4$. Khẳng định nào sau đây đúng?
}{%
\begin{tasks}(2)
    \task $(S_1)$ và $(S_2)$ cắt nhau
    \task $(S_2)$ nằm trong lòng $(S_1)$ (đựng nhau)
    \task $(S_1)$ và $(S_2)$ tiếp xúc trong nhau
    \task $(S_1)$ và $(S_2)$ ở ngoài nhau
\end{tasks}
}
\drawdashlinesfull

% =========================================================================
% BẢNG ĐÁP ÁN BTDD PHƯƠNG TRÌNH MẶT CẦU (ĐÃ KIỂM TRA ĐỘC LẬP - ẨN TRONG CODE)
% -------------------------------------------------------------------------
% 1.B   2.A   3.A   4.A   5.A   6.B   7.C   8.A   9.B   10.A
% 11.A  12.B  13.B  14.C  15.C  16.B  17.B  18.B  19.C  20.B
% =========================================================================

\end{document}